\section{Mathematical Foundations}
\label{sec:foundations}

\subsection{Fractional Brownian motion and its covariance structure}

Fractional Brownian motion $W_t^H$ is the unique (up to scaling) continuous Gaussian
process with stationary increments and the self-similarity property
$W_{ct}^H \stackrel{d}{=} c^H W_t^H$ for all $c > 0$
\citep{mandelbrot1968fractional}.  Its covariance kernel is
\begin{equation}
  \label{eq:fbm-cov}
  \mathbb{E}[W_t^H W_s^H]
  = \tfrac{1}{2}\!\left(t^{2H} + s^{2H} - |t-s|^{2H}\right),
\end{equation}
which is symmetric, positive-definite for $H \in (0,1)$, and reduces to the standard
Brownian covariance $\min(t,s)$ at $H = \nicefrac{1}{2}$.  The self-similarity
$W_{ct}^H \stackrel{d}{=} c^H W_t^H$ and the stationarity of increments
$W_{t+h}^H - W_t^H \stackrel{d}{=} W_h^H$ are the properties that make fBM a
tractable model of rough volatility; §\ref{sec:toeplitz} exploits increment
stationarity directly.  The critical property for simulation is that fBM is
non-Markovian for $H \neq \nicefrac{1}{2}$: the value of $W_t^H$ depends on
the entire history of the driving noise, not just its most recent value.  This rules
out any forward recursion and forces each sample path to be drawn from a joint
$N$-dimensional Gaussian.

Discretizing $[0,T]$ into $N$ equally spaced points $t_j = jT/N$ yields an $N\times N$
covariance matrix $C$ with $C_{jk} = \mathbb{E}[W_{t_j}^H W_{t_k}^H]$ as in
equation~\eqref{eq:fbm-cov}.  The matrix $C$ is dense (every pair of time points is correlated) and positive-definite
for all $H \in (0,1)$ and $N \geq 1$.  This density explains the $O(N^2)$ per-path
cost: sampling from
$\mathcal{N}(0,\nu^2 C)$ by any exact method requires touching all $N^2/2$ entries
of $C$.  Algorithm~1 (§\ref{sec:cholesky}) accepts this cost and factorizes $C$
exactly; Algorithms~2 and~3 avoid it by exploiting structure that $C$ possesses.

\subsection{Stationarity of fGn and the Toeplitz property}
\label{sec:toeplitz}

Although the fBM covariance matrix is dense, its increment process has a
crucial structural simplification.  Define the fractional Gaussian noise (fGn)
increments $\delta W_m = W_{(m+1)\Delta t}^H - W_{m\Delta t}^H$ for $m = 0, \ldots,
N-1$, where $\Delta t = T/N$.  Stationarity of fBM increments implies
\begin{equation}
  \label{eq:fgn-cov}
  \gamma(k) \coloneqq \mathrm{Cov}(\delta W_m,\, \delta W_{m+k})
  = \frac{\Delta t^{2H}}{2}\bigl(|k+1|^{2H} + |k-1|^{2H} - 2|k|^{2H}\bigr),
\end{equation}
which depends only on the lag $k = |m - n|$, not on the absolute index $m$
\citep{wood1994simulation}.  The $N\times N$ covariance matrix of
$(\delta W_0, \ldots, \delta W_{N-1})$ is therefore Toeplitz: its $(m,n)$
entry equals $\gamma(|m-n|)$, so every diagonal is constant.

Figure~\ref{fig:structure} illustrates the contrast.  Panels~(a) and~(b) show the
$N\times N$ covariance heatmaps of fBM and fGn respectively.  The fBM matrix is not
Toeplitz: row $m$ of $C$ varies with $m$ because the variance $t_m^{2H}$ grows with
time.  The fGn matrix has constant diagonals, confirming the Toeplitz property exactly.
Panel~(c) shows diagonal slices $C[m, m+k]$ vs.\ lag $k$ for two different rows
$m$: the fGn curves overlap to numerical precision, while the fBM curves diverge.

\begin{figure}[t]
  \centering
  \includegraphics[width=\linewidth]{structure_analysis.png}
  \caption{Structural properties of the fBM and fGn covariance matrices ($H=0.10$).
    \textbf{(a)} fBM covariance $C(t_m, t_n)$: dense, not Toeplitz (variance grows
    with time).
    \textbf{(b)} fGn covariance $\gamma(|m-n|)$: Toeplitz (constant diagonals).
    \textbf{(c)} Diagonal slice $C[m, m+k]$ vs.\ lag $k$: fGn rows for $m=0$ and
    $m=N/4$ overlap (Toeplitz); fBM rows diverge (non-stationary).
    \textbf{(d)} Normalised singular values $\sigma_k/\sigma_1$ of the off-diagonal
    block (solid) and normalised eigenvalues of the full matrix $C$ (dashed) at
    $H=0.10$ and $H=0.50$.  The off-diagonal block compresses to rank $\leq 3$ at
    both $H$; the full spectrum decays slowly, most slowly at $H=0.10$.
    Controlled: $H=0.10$, $N=64$ (heatmaps), $N=128$ (SVD).
    Independent: process type or Hurst exponent.
    Dependent: covariance pattern and singular value magnitude.}
  \label{fig:structure}
\end{figure}

Algorithm~2 (§\ref{sec:fft}) depends on this Toeplitz structure: only Toeplitz
matrices embed into a circulant diagonalizable by the discrete Fourier transform.
Applying the FFT method directly to the fBM covariance matrix 
produces negative eigenvalues and pathological path output.

\subsection{Circulant embedding and its positivity condition}
\label{sec:circulant}

A Toeplitz matrix $T$ with first row $(\gamma(0), \gamma(1), \ldots, \gamma(N-1))$
can be embedded into a $2N\times 2N$ circulant matrix $C_\mathrm{circ}$ whose
first row is the symmetric extension
\begin{equation}
  c = \bigl(\gamma(0),\, \gamma(1),\, \ldots,\, \gamma(N-1),\, 0,\,
             \gamma(N-1),\, \ldots,\, \gamma(1)\bigr).
\end{equation}
Circulant matrices are diagonalized by the DFT: the eigenvalues of $C_\mathrm{circ}$
are exactly $\lambda = \mathrm{FFT}(c)$, computable in $O(N\log N)$ operations.
To sample, draw $W = a + ib$ with $a, b \sim \mathcal{N}(0, I_{2N})$ independent
(so each component has real and imaginary parts that are independent standard
normals),
scale component-wise by $\sqrt{\lambda_j / (2N)}$, apply the inverse FFT, and take
the real part of the first $N$ elements \citep{davies1987tests}.  One forward FFT (once per experiment) to compute
$\{\lambda_j\}$, then one inverse FFT per path: total $O(N\log N)$ per path.

The method is valid only if every eigenvalue $\lambda_j$ is non-negative, so that
$\sqrt{\lambda_j}$ is real \citep{wood1994simulation}.  For rough $H$ this holds at
every $N$, not just asymptotically.  \citet{craigmile2003simulating} proves that the
Davies--Harte embedding is non-negative definite for any stationary process whose
autocovariance is non-positive at all non-zero lags, a class that includes fGn with
$H < \nicefrac{1}{2}$.  The argument is short enough to give in full for the
embedding used here.  Because $x \mapsto x^{2H}$ is concave for $H < \nicefrac{1}{2}$,
the second difference in equation~\eqref{eq:fgn-cov} is negative, so $\gamma(k) < 0$
for every $k \geq 1$.  The eigenvalues of the symmetric circulant are
\begin{equation}
  \lambda_j = \gamma(0) + 2\sum_{k=1}^{N-1} \gamma(k)\cos\!\left(\frac{\pi jk}{N}\right),
  \qquad j = 0, \ldots, 2N-1,
\end{equation}
and because each $\gamma(k)$ is negative, replacing every cosine by $1$ can only lower
the sum.  The smallest eigenvalue is therefore $\lambda_0$, and the sum telescopes:
\begin{equation}
  \label{eq:lambda-min}
  \begin{split}
  \lambda_{\min} = \lambda_0
    &= \gamma(0) + 2\sum_{k=1}^{N-1}\gamma(k) \\
    &= \Delta t^{2H}\bigl(N^{2H} - (N-1)^{2H}\bigr)
     \approx 2H\,\Delta t^{2H} N^{2H-1} > 0.
  \end{split}
\end{equation}
The FFT method is thus exact on the grid for every $H \leq \nicefrac{1}{2}$ (at
$H = \nicefrac{1}{2}$ the increments are white noise and every $\lambda_j = \gamma(0)$).
The C++ pricer checks the sign of every eigenvalue and throws rather than clipping.

Equation~\eqref{eq:lambda-min} also shows that the margin is thin.  Relative to the
increment variance $\gamma(0) = \Delta t^{2H}$, the smallest eigenvalue is
$N^{2H} - (N-1)^{2H}$: $0.24\%$ at $H = 0.10$, $N = 252$ and $0.08\%$ at $N = 1000$,
decaying like $N^{-(1-2H)}$.  This mirrors the fGn spectral density, which vanishes
at zero frequency when $H < \nicefrac{1}{2}$.  Section~\ref{sec:stability} confirms
the formula numerically.  The key pitfall is $\gamma(0)$ itself: dropping the
$|k-1|^{2H}$ term at $k = 0$ halves it, which breaks the bound and spuriously makes
$28\%$ of the eigenvalues negative at $H = 0.10$.  An earlier version of my
stability analysis made exactly this mistake and concluded, wrongly, that the FFT
method needed eigenvalue clipping.

\subsection{Low-rank structure of smooth off-diagonal blocks}
\label{sec:lowrank}

The fBM kernel $C(t,s) = \tfrac{1}{2}(t^{2H}+s^{2H}-|t-s|^{2H})$ has a singularity at $t = s$, 
since the term $|t-s|^{2H}$ introduces one. For $H = 0.10$ this gives $|t-s|^{0.2}$, whose derivative
$0.2|t-s|^{-0.8}$ is unbounded as $t \to s$. The function is H\"older-continuous
with exponent $0.2$ but strictly non-differentiable on the diagonal.

Off-diagonal blocks of $C$ (submatrices coupling two disjoint time intervals) lie
in the smooth region $t \neq s$.  Smooth bivariate kernels generate matrices with
rapidly decaying singular values on such blocks \citep{candes2009fast}.

Figure~\ref{fig:structure}(d) confirms this, and locates the difficulty precisely.
The off-diagonal block coupling the first half of the grid to the second half is
exactly rank~1 at $H = 0.50$, where the kernel is $\min(t,s)$.  At $H = 0.10$ its
singular values still fall geometrically: below $1\%$ of $\sigma_1$ by rank $k = 3$
and below $10^{-4}$ by $k = 5$.  The spectrum of the full matrix $C$ behaves very
differently.  It decays only algebraically, and more slowly at $H = 0.10$ than at
$H = 0.50$.  The slow decay that limits a global low-rank approximation therefore
comes from the near-diagonal singularity, not from the far field.  This is exactly
the structure a hierarchical matrix exploits (§\ref{sec:hmatrix-future}).